% SPDX-License-Identifier: CC-BY-SA-4.0
% Author: Matthieu Perrin

\begingroup

\begin{exercice}[Equipotence]\label{exo:introduction/maths/equipotence}

  Soient $E$ un ensemble et $A, B \in \mathscr{P}(E)$ deux sous-ensembles de $E$. On dit que \emph{$A$ est équipotent à $B$}, noté
  $A \sim B$, s'il existe une bijection de $A$ dans $B$.
  \begin{question}
  \item Lister tous les éléments de $\mathscr{P}(\{a, b, c\})$.
  \item Si $E = \{a, b, c\}$, montrer que $\{a, b\} \sim \{b, c\}$. 
  \item Montrer que, pour tout ensemble $E$, $\sim$ est une relation d'équivalence sur  $\mathscr{P}(E)$. 
  \item Déterminer l'ensemble quotient de $\mathscr{P}(\{a, b, c\})$ par $\sim$. 
  \end{question}

\end{exercice}

\begin{correction}

  \begin{question}
  \item $\emptyset, \{a\}, \{b\}, \{c\}, \{a, b\}, \{a, c\}, \{b, c\}, \{a, b, c\}$
  \item La fonction $\left\{\begin{array}{rcl}
        a & \longmapsto & b\\
        b & \longmapsto & c
      \end{array}\right\}$ est une bijection de $\{a, b\}$ vers $\{b, c\}$, donc $\{a, b\} \sim \{b, c\}$.
  \item Il faut montrer les trois propriétés des relations d'équivalence :
    \begin{description}
    \item[Réflexivité :] Soit $A \in \mathscr{P}(E)$.  
      La fonction identité $\left\{\begin{array}{rcl}
        A & \longmapsto & A\\
        x & \longmapsto & x
      \end{array}\right\}$ est bijective, donc $A \sim A$.
    \item[Symétrie :] Soient $A, B \in \mathscr{P}(E)$ tels que $A \sim B$.  
      Il existe une bijection $f$ de $A$ dans $B$. Sa réciproque $f^{-1}$ est une bijection de $B$ dans $A$, donc $B \sim A$.
    \item[Transitivité :] Soient $A, B, C \in \mathscr{P}(E)$ tels que $A \sim B$ et $B \sim C$.  
      Il existe deux bijections $f : A \to B$ et $g : B \to C$.  
      Posons $h = g\circ f = \left\{\begin{array}{rcl}
        A & \longmapsto & C \\
        x & \longmapsto & g(f(x))
      \end{array}\right\}$.  
      $h$ est bijective, donc $A \sim C$.
    \end{description}
  \item $\big\{\{\emptyset\},\ \{\{a\}, \{b\}, \{c\}\},\ \{\{a, b\}, \{a, c\}, \{b, c\}\},\ \{\{a, b, c\}\}\big\}$
  \end{question}

\end{correction}

\endgroup
\endinput
